\ifmostrarGabarito
\textbf{Resposta: (B)}

\textbf{Justificativa:}
No bloco \code{if/elif/else}, a execução ocorre de cima para baixo e para no \textbf{primeiro} ramo cuja condição for verdadeira. Como \code{pontos = 75}, a primeira condição (\code{pontos >= 90}) é falsa, mas a segunda (\code{pontos >= 70}) é verdadeira. O valor \code{"Prata"} é atribuído e o restante da estrutura é ignorado.

\fi
